\textbf{\textit{Soal. }}For $\triangle ABC$ and its circumcircle $\omega$, draw the tangents at $B,C$ to $\omega$ meeting at $D$. Let the line $AD$ meet the circle with center $D$ and radius $DB$ at $E$ inside $\triangle ABC$. Let $F$ be the point on the extension of $EB$ and $G$ be the point on the segment $EC$ such that $\angle AFB=\angle AGE=\angle A$. Prove that the tangent at $A$ to the circumcircle of $\triangle AFG$ is parallel to $BC$.


\begin{proof}[\textbf{Solusi.}] (Sunday, 16 October 2022) Misalkan $M = CE \cap AF$ dan $N = BE \cap AG$.\\
Perhatikan bahwa $\angle DBC = \angle BCD = \angle BAC$. Sehingga didapat $\angle BEC = 180^\circ-\dfrac{1}{2}\angle BDC = 90^\circ+\dfrac{1}{2}(180^\circ-\angle BDC) = 90^\circ+\dfrac{1}{2}(2\angle DBC) = 90^\circ+\angle DBC = 90^\circ+\angle BAC$ yang menyebabkan $\angle GNF = 180^\circ -\angle GEN -\angle AGE = \angle BEC - \angle BAC = 90^\circ$. Dengan cara serupa didapat $\angle GMF = 90^\circ$. Oleh karena itu, $G,N,M,F$ konsiklis. 

Selanjutnya, karena $\angle EMA = 90^\circ$ maka $\angle MNB = \angle MNE = \angle MAE = 90^\circ - \angle AEM = 90^\circ -\angle DEC$. Padahal, karena $D$ pusat lingkaran $(BEC)$ maka $\angle NBC = \angle EBC = \dfrac{1}{2}\angle EDC = 90^\circ - \angle DEC$. Ini berarti $\angle MNB = \angle NBC$. Didapatkan bahwa $MN \parallel BC$. Oleh karena itu, cukup dibuktikan bahwa garis singgung $(AFG)$ di $A$ sejajar dengan $MN$.

\begin{center}
    \begin{asy}
          /* Geogebra to Asymptote conversion, documentation at artofproblemsolving.com/Wiki go to User:Azjps/geogebra */
        import graph; size(10cm); 
        real labelscalefactor = 0.5; /* changes label-to-point distance */
        pen dps = linewidth(0.7) + fontsize(10); defaultpen(dps); /* default pen style */ 
        pen dotstyle = black; /* point style */ 
        real xmin = -9, xmax = 6, ymin = -7.558813578242471, ymax = 10;  /* image dimensions */
        pen zzttqq = rgb(0.6,0.2,0); pen zzttff = rgb(0.6,0.2,1); 
        
        draw((-2.882744078244634,7.835999567379222)--(-3.2525786115774866,0.2686160391839343)--(3.063057265337379,0.21171841867118776)--cycle, linewidth(1) + zzttqq); 
         /* draw figures */
        draw((-2.882744078244634,7.835999567379222)--(-3.2525786115774866,0.2686160391839343), linewidth(1) + zzttqq); 
        draw((-3.2525786115774866,0.2686160391839343)--(3.063057265337379,0.21171841867118776), linewidth(1) + zzttqq); 
        draw((3.063057265337379,0.21171841867118776)--(-2.882744078244634,7.835999567379222), linewidth(1) + zzttqq); 
        draw(circle((-0.061740540954276925,3.905401899328803), 4.838146234405085), linewidth(0.4) + linetype("4 4")); 
        draw((-2.882744078244634,7.835999567379222)--(-0.11927101732279266,-2.4804809775764753), linewidth(1)); 
        draw(circle((-0.11927101732279266,-2.4804809775764753), 4.168351099388716), linewidth(0.4) + linetype("2 2")); 
        draw((-3.2525786115774866,0.2686160391839343)--(-0.11927101732279266,-2.4804809775764753), linewidth(1)); 
        draw((3.063057265337379,0.21171841867118776)--(-0.11927101732279266,-2.4804809775764753), linewidth(1)); 
        draw((-1.197821224442296,1.5459169583580241)--(3.063057265337379,0.21171841867118776), linewidth(1)); 
        draw((-2.882744078244634,7.835999567379222)--(1.5651849928587422,0.6807435277823698), linewidth(1)); 
        draw((-2.882744078244634,7.835999567379222)--(-5.735625030892791,-1.2749226255755657), linewidth(1)); 
        draw((1.5651849928587422,0.6807435277823698)--(-5.735625030892791,-1.2749226255755657), linewidth(1)); 
        draw((-4.526016563571792,2.5880666868495124)--(-1.197821224442296,1.5459169583580241), linewidth(1)); 
        draw((-1.197821224442296,1.5459169583580241)--(0.407160934825726,2.543623646323409), linewidth(1)); 
        draw((-1.197821224442296,1.5459169583580241)--(-5.735625030892791,-1.2749226255755657), linewidth(1)); 
        draw((-4.526016563571792,2.5880666868495124)--(0.407160934825726,2.543623646323409), linewidth(1));
        draw(circle((-2.0852200190170245,-0.29708954889660094), 3.77910231094306), linewidth(0.8)); 
        draw(circle((-2.9276814459181923,2.847951755614001), 4.988250228233267), linewidth(1.2) + linetype("0 3 4 3") + zzttff); 
        draw((xmin, -0.009009009009008551*xmin + 7.81002890000765)--(xmax, -0.009009009009008551*xmax + 7.81002890000765), linewidth(1)); /* line */
         /* dots and labels */
        dot((-2.882744078244634,7.835999567379222),dotstyle); 
        label("$A$", (-2.7961641842937297,8.077592575921365), NE * labelscalefactor); 
        dot((-3.2525786115774866,0.2686160391839343),dotstyle); 
        label("$B$", (-3.1651649194952367,0.4900149583403892), NW * labelscalefactor); 
        dot((3.063057265337379,0.21171841867118776),dotstyle); 
        label("$C$", (3.1539726708305733,0.4438898664402009), SE * labelscalefactor); 
        dot((-0.11927101732279266,-2.4804809775764753),linewidth(4pt) + dotstyle); 
        label("$D$", (-0.028658670282425858,-2.3005531016210035), SE * labelscalefactor); 
        dot((-1.197821224442296,1.5459169583580241),linewidth(4pt) + dotstyle); 
        label("$E$", (-1.1125983299368531,1.7353924396454734), NE * labelscalefactor); 
        dot((1.5651849928587422,0.6807435277823698),linewidth(4pt) + dotstyle); 
        label("$G$", (1.6549071840744507,0.8590156935418957), NE * labelscalefactor); 
        dot((-5.735625030892791,-1.2749226255755657),linewidth(4pt) + dotstyle); 
        label("$F$", (-5.632857336155316,-1.1013007122161074), NE * labelscalefactor); 
        dot((-4.526016563571792,2.5880666868495124),linewidth(4pt) + dotstyle); 
        label("$M$", (-4.4336049467504175,2.7732070073997104), NW * labelscalefactor); 
        dot((0.407160934825726,2.543623646323409),linewidth(4pt) + dotstyle); 
        label("$N$", (0.5017798865697407,2.727081915499522), NE * labelscalefactor); 
        clip((xmin,ymin)--(xmin,ymax)--(xmax,ymax)--(xmax,ymin)--cycle); 
         /* end of picture */
    \end{asy}
\end{center}

Sekarang, jika dimisalkan $l$ adalah garis singgung lingkaran $(AGF)$ di $A$. Dari Tangent-Secant Theorem dan dari fakta $M,N,G,F$ konsiklis, $\angle (GA, l)=\angle GFA = \angle GFM = 180^\circ - \angle MNG = \angle ANM$ yang menunjukkan $l \parallel MN$. Terbukti.

    
\end{proof}